\chapter{The \emph{Chuci} and Qu Yuan}

\section{The southern tradition}

If the \emph{Shijing} is the poetry of the northern Zhou heartland, the \zht{楚辭}
\emph{Chǔcí} (\emph{Songs of Chu}, \emph{Elegies of Chu}) is the poetry of the southern
state of Chu \zht{楚}, in the middle Yangzi region, compiled in the Han dynasty (the
standard edition is by Wang Yi \zht{王逸}, 2nd century CE). Chu culture was regarded by
northerners as exotic: shamanism, river goddesses, lush flora, a different dialect.

\begin{center}
\begin{tabular}{lll}
\toprule
 & \emph{Shijing} \zht{詩經} & \emph{Chuci} \zht{楚辭} \\
\midrule
Region & North (Yellow River) & South (Yangzi, Chu) \\
Date & c.\ 1000--600 BCE & 4th--3rd c.\ BCE (and Han imitations) \\
Authorship & Anonymous & Attributed to Qu Yuan and followers \\
Line & Four characters, regular & Longer, variable, with \zht{兮} \emph{xī} \\
Voice & Communal, social & Individual, lyrical, visionary \\
Imagery & Farm, court, marriage & Fragrant plants, shamans, sky journeys \\
\bottomrule
\end{tabular}
\end{center}

\begin{important}[The two sources]
The \emph{Shijing} and the \emph{Chuci} are the twin sources of the Chinese poetic
tradition: the \emph{Shijing} for the plain, morally engaged lyric; the \emph{Chuci} for
the impassioned, imaginative, self-dramatising voice. Later critics use the shorthand
\zht{風騷} \emph{fēng sāo} (``\emph{Airs} and \emph{Li Sao}'') for poetry itself.
\end{important}

\section{Qu Yuan \zht{屈原} (c.\ 340--278 BCE)}

The story of Qu Yuan comes mainly from Sima Qian's \emph{Records of the Grand Historian}
(\zht{史記}, c.\ 100 BCE), written two centuries later:
\begin{enumerate}
  \item Qu Yuan, a noble of Chu, served King Huai \zht{楚懷王} loyally and advised against
        trusting the state of Qin.
  \item Slandered by jealous courtiers, he was dismissed and later exiled to the south.
  \item King Huai ignored his advice, went to Qin and died a captive; Qin eventually
        sacked the Chu capital (278 BCE).
  \item In despair, Qu Yuan clasped a stone and drowned himself in the Miluo River
        \zht{汨羅江}.
\end{enumerate}

\begin{aside}
The Dragon Boat Festival (\zht{端午節}, fifth day of the fifth month) is traditionally
said to commemorate the search for Qu Yuan's body: boats raced to find him, and people
threw rice dumplings (\zht{粽子} \emph{zòngzi}) into the river so the fish would eat them
instead of the poet. Whether or not the festival really began this way, it made Qu Yuan
the patron saint of the loyal, misunderstood official.
\end{aside}

\begin{remark}
Modern scholars debate how many \emph{Chuci} poems Qu Yuan actually wrote, and some have
doubted he existed at all. For this course, ``Qu Yuan'' is above all a \emph{persona}: the
upright minister rejected by his ruler, which later poets (Jia Yi, Li Bai, Du Fu) adopt
when they are themselves exiled.
\end{remark}

\section{\emph{Li Sao} \zht{離騷} ``Encountering Sorrow''}

The \emph{Li Sao} is the longest poem of pre-Qin China (373 lines). The title is
already disputed: \zht{離} ``to meet with'' or ``to leave'', \zht{騷} ``sorrow,
trouble'', hence ``Encountering Sorrow'' or ``Leaving Sorrow behind''.

\subsection{Opening: a genealogy}

\begin{center}
\begin{tabular}{p{0.38\textwidth}p{0.52\textwidth}}
\toprule
\zht{帝高陽之苗裔兮，} & A distant descendant of the god-king Gaoyang, \\
\zht{朕皇考曰伯庸。} & my late honoured father was called Boyong. \\
\zht{攝提貞于孟陬兮，} & When the year-star stood in the first month of spring, \\
\zht{惟庚寅吾以降。} & on the day \emph{gengyin}, I came down to earth. \\
\bottomrule
\end{tabular}
\end{center}

The speaker begins with his divine ancestry and auspicious birth, then tells how he was
given fine names and adorned himself with fragrant plants.

\subsection{The line and the particle \zht{兮}}

The typical \emph{Sao}-style couplet has a long first line ending in the breath particle
\zht{兮} \emph{xī} (roughly ``ah!'', a sung pause) and a second line without it:

\begin{center}
\begin{tikzpicture}[ch/.style={draw, minimum size=0.7cm}, xi/.style={draw, fill=orange!30, minimum size=0.7cm},
  gl/.style={font=\scriptsize}]
  \foreach \c [count=\i] in {朝,搴,阰,之,木,蘭} { \node[ch] at (0.75*\i,1.2) {\zht{\c}}; }
  \node[xi] at (0.75*7,1.2) {\zht{兮}};
  \foreach \c [count=\i] in {夕,攬,洲,之,宿,莽} { \node[ch] at (0.75*\i,0) {\zht{\c}}; }
  \draw[thick] (0.4,1.65) -- (2.6,1.65) node[midway, above, gl] {when, verb, where};
  \draw[thick] (2.65,1.65) -- (4.85,1.65) node[midway, above, gl] {\zht{之} + object};
  \node[gl, anchor=west] at (5.6,1.6) {breath particle};
  \node[gl, anchor=west] at (6.2,0.8) {6 + \zht{兮} / 6};
  \node[gl, anchor=north west] at (0.35,-0.5) {dusk $\cdot$ gather $\cdot$ isle $\cdot$ 's $\cdot$ winter-grass};
\end{tikzpicture}
\end{center}

\begin{center}
\begin{tabular}{p{0.38\textwidth}p{0.52\textwidth}}
\toprule
\zht{汨余若將不及兮，} & Swiftly I hurried, as if I could not keep up, \\
\zht{恐年歲之不吾與。} & afraid the years would not wait for me. \\
\zht{朝搴阰之木蘭兮，} & At dawn I plucked magnolia on the hill, \\
\zht{夕攬洲之宿莽。} & at dusk I gathered winter-grass on the isle. \\
\bottomrule
\end{tabular}
\end{center}

Note the morning/evening (\zht{朝} / \zht{夕}) parallelism, which the \emph{Chuci}
uses constantly, and the anxiety about time passing.

\subsection{The system of fragrant plants}

\begin{important}[Symbolism in the \emph{Li Sao}]
\begin{itemize}
  \item \textbf{Fragrant plants} (orchid \zht{蘭}, angelica \zht{芷}, magnolia
        \zht{木蘭}): the speaker's virtue. To wear or plant them is to cultivate one's
        moral nature.
  \item \textbf{Weeds and foul plants}: the slanderers at court.
  \item \textbf{The Fair One} (\zht{美人}): the ruler, or an ideal partner. The speaker
        often speaks as a rejected lover or wife; the political relationship of minister
        and ruler is cast as courtship. (This lover/ruler allegory becomes a permanent
        convention.)
  \item \textbf{The sky journey}: rejected on earth, the speaker drives a chariot drawn by
        dragons and phoenixes through the heavens, seeking a goddess or a worthy match,
        and fails each time.
\end{itemize}
\end{important}

\begin{center}
\begin{tabular}{p{0.38\textwidth}p{0.52\textwidth}}
\toprule
\zht{長太息以掩涕兮，} & Long I sighed and wiped away my tears, \\
\zht{哀民生之多艱。} & grieving that the people's life is so hard. \\
\bottomrule
\end{tabular}
\end{center}

\subsection{Structure and ending}

\begin{enumerate}
  \item Self-presentation: ancestry, birth, fragrant adornment.
  \item Political lament: the ruler is misled; the speaker refuses to compromise.
  \item Consultations: a sister figure (\zht{女嬃}) advises him to bend; he appeals to the
        ancient sage-king Shun.
  \item Cosmic journey: through the heavens, seeking goddesses; the gate of heaven is shut.
  \item Divination: a shaman tells him to leave Chu for a better lord.
  \item Departure and return: as he rises in splendour, he glimpses his old homeland, and
        his servant and horses grieve and will not go on.
  \item Envoi (\zht{亂曰}): ``Enough! There is no one in the state who knows me \ldots\ I
        will follow Peng Xian to his dwelling.'' (Peng Xian, a legendary worthy said to
        have drowned himself; read as foreshadowing Qu Yuan's suicide.)
\end{enumerate}

\section{The \emph{Nine Songs} \zht{九歌}}

The \emph{Nine Songs} (actually eleven pieces) are shorter shamanistic hymns to Chu
deities: the Lord of the East (the sun), the Lord of the Clouds, the Lady of the Xiang
River, the Mountain Spirit and others. Many seem to be performed by a shaman who invites the
god, enjoys a brief ecstatic union and then is left alone. The pattern of longing for a
divine lover who does not come is the emotional template of the \emph{Li Sao}.

\begin{example}[A characteristic line from ``Lady of the Xiang'' \zht{湘夫人}]
\zht{嫋嫋兮秋風，洞庭波兮木葉下。} \\
``Gently, gently, the autumn wind; / on Dongting's waves the leaves of the trees fall.''
Here \zht{兮} sits \emph{inside} the line. This couplet founded the image of autumn
leaves on Lake Dongting, endlessly echoed by Tang poets.
\end{example}

\section{From \emph{Chuci} to \emph{fu}}

The Han dynasty imitated the \emph{Chuci} extensively and developed from it the
\zht{賦} \emph{fù} ``rhapsody'' (the same character as the \emph{fu} of fu-bi-xing):
long, elaborately descriptive compositions mixing prose and verse, cataloguing palaces,
hunting parks and capital cities. Sima Xiangru \zht{司馬相如} (179--117 BCE) is the master.
We will not read Han rhapsodies in detail, but you should know the line of descent:
\emph{Chuci} $\to$ Han \emph{fu}.

\section*{Summary}
\begin{itemize}
  \item \emph{Chuci}: southern (Chu) songs compiled in the Han; long lines with \zht{兮};
        individual, visionary voice.
  \item Qu Yuan: loyal minister slandered and exiled, drowned in the Miluo; model for all
        later frustrated officials; Dragon Boat Festival.
  \item \emph{Li Sao}: fragrant plants = virtue; Fair One = ruler; sky journey; final
        decision to follow Peng Xian.
  \item \emph{Nine Songs}: shamanic hymns, longing for an absent deity.
  \item Pairing: \emph{Shijing} (north, communal) and \emph{Chuci} (south, individual):
        \zht{風騷}.
\end{itemize}
